\section{Method}
\label{sec:method}
\subsection{Problems, instances and hosts}
\label{sec:instances}
In $P||\Cmax$, $n$ jobs with processing times $p_j$ are assigned to $m$ identical machines to minimise $\Cmax=\max_i\sum_{j\in S_i}p_j$; the classical bound is $\mathrm{LB}=\max(\max_jp_j,\lceil\sum_jp_j/m\rceil)$~\cite{Graham1969}. The \emph{synthetic} instances have $p_j\sim U[1,100]$ in twelve classes $(n,m)$ from $(20,4)$ to $(300,20)$, with $100$ test and $30$ separate tuning instances per class; the \emph{published} families B, C, E and F of the benchmark reviewed by \citet{OstojicEtAl2025} give disjoint sets of $525$ tuning, $2{,}550$ discovery and $3{,}950$ held-out test instances, $5{,}720$ instances on which an offline selector is trained, and $780$ of family F. In $R||\Cmax$ the time $p_{jk}$ depends on the machine: we use (a) the twelve synthetic classes with $p_{jk}\sim U[1,100]$, (b) three published families (uncorrelated, correlated jobs, correlated machines~\citep{AbladEtAl2021}; $25$ sizes each, generated by us from the description), and (c) the \emph{standard benchmark} of Fanjul-Peyro and Ruiz~\citep{FanjulPeyroRuiz2010}: $1{,}400$ instances of seven types with $m\in\{10,\dots,50\}$ and $n\in\{100,\dots,1000\}$, with best known values from CPLEX, Gurobi and an exact Lagrangian method run by third parties (the minimum of the three is proven optimal for $73\%$ of the test instances); two instances per cell tune and five test ($280$ tuning and $700$ test instances). The other hosts are a permutation flowshop (nine synthetic classes and the $90$ instances of \citet{Taillard1993}), a steady-state genetic algorithm for $P||\Cmax$ ($1{,}200$ instances) and, outside scheduling, an iterated local search for the Euclidean TSP on generated instances and on $37$ instances of the TSPLIB (Section~\ref{sec:tsp}).

\subsection{Search mechanisms and learners}
\label{sec:learners}
The \textbf{IG} starts from LPT and alternates destruction (removing $d$ jobs) and greedy reconstruction, accepting a candidate that does not worsen the makespan. The \textbf{ILS} chooses at each iteration one of four neighbourhood operators. The \emph{actions} of the IG are the destruction strategy ($3$), the destruction size ($5$), their $15$ combinations and these with or without a descent after reconstruction ($30$); on the standard benchmark the sizes scale with $n$ (B5, B15, B30). At each iteration a selector chooses an action from a state of $81$ values (gap to the bound, load imbalance, fraction of the budget used, stagnation), receives as reward the improvement of the makespan truncated at zero, and the selector is the only component that changes between arms. The learners are tabular Q-Learning ($\alpha=0.1$, $\gamma=0.9$, $\varepsilon$ from $1.0$ to $0.05$)~\cite{WatkinsDayan1992}, UCB1~\cite{AuerEtal2002}, a policy gradient~\cite{Williams1992}, six stateless methods (probability matching, adaptive pursuit, an ALNS roulette, discounted Thompson sampling~\cite{Russo2018}, DMAB~\cite{DaCostaEtal2008}, online racing), LinUCB~\cite{LiEtAl2010}, a DQN~\cite{MnihEtAl2015}, an optimistic Q-Learning (optimistic table of initial value $c$, constant $\varepsilon=0.05$, scale-invariant reward) and two re-implemented published designs (an episodic Q-Learning in the style of \citet{KarimiMamaghanEtal2023} and an FRRMAB-style bandit~\cite{LiFialhoEtal2014}); hyperparameters are tuned on tuning instances and then frozen. \emph{Random choice} over the same actions is the control; the \emph{static competitor} is the best fixed action per class on the tuning instances.

\paragraph{Action order and ties} The standard Q-Learning starts from a table of zeros, its reward is non-negative and its greedy step takes the \emph{first} maximum; in an unvisited state all values are equal, so after the exploration rate has decayed it repeats the action with the smallest index (the smallest and cheapest, as our lists are ordered by size). We therefore report the standard implementation (\textsc{ql}) and an order-invariant version (\textsc{ql-perm}) that permutes the actions at random in each run, which is equivalent to breaking ties at random.

\subsection{The budget-fair attribution protocol}
\label{sec:protocol}
A \emph{class} is an $(n,m)$ pair or a published family by size tercile; a \emph{cell} is a class in one action space; the \emph{validated action} is the best fixed action per class, chosen on separate instances; the \emph{headroom} of a cell is the gain over random choice of the best fixed action per instance. A \emph{gain} is $-\bar d$ in percentage points of the makespan, positive when the first-named arm is better; \emph{registered} results were specified before the battery was run, \emph{post hoc} ones were not.
\begin{enumerate}
\item \textbf{Random control at a matched budget.} The effect of the policy is the difference between the learner and uniform random choice over the same actions, iterations and instances.
\item \textbf{Equal CPU time and the decision/cost split.} The ratio $K$ of the CPU time of $1000$ iterations of one arm to that of another is measured on separate instances and the second arm receives $\mathrm{round}(1000K)$ iterations. With $\bar d_{\mathrm{iter}}$ and $\bar d_{\mathrm{CPU}}$ the mean relative differences of the learner against random choice at equal iterations and at equal CPU time,
\begin{equation}
G_{\mathrm{CPU}}=G_{\mathrm{dec}}+G_{\mathrm{cost}},\qquad G_{\mathrm{dec}}=-\bar d_{\mathrm{iter}},\qquad G_{\mathrm{cost}}=\bar d_{\mathrm{iter}}-\bar d_{\mathrm{CPU}}.
\label{eq:decomp}
\end{equation}
A learner with $G_{\mathrm{dec}}\le0$ and $G_{\mathrm{cost}}>0$ wins at equal CPU time without having decided anything better.
\item \textbf{A static competitor validated offline}: the best fixed action per class and, when actions are complementary, a uniform mixture of the best few; at equal CPU time it receives its own cost ratio.
\item \textbf{Cross-fitted oracle headroom.} The action is chosen on two seeds and evaluated on two fresh ones (Section~\ref{sec:estimator}).
\item \textbf{Power and positive controls}: the minimum detectable gain, analytic and by planting gains in the paired differences, and a planted reward bonus.
\item \textbf{Order sensitivity}: the share of choices on the action listed first, and the learner with the actions permuted at random (Section~\ref{sec:order}).
\item \textbf{What the static competitor knows}: static policies chosen with decreasing knowledge of the instance (one action for all, one per size, per type, per class), in transfer regimes (the tested type or size withheld), compared with the learner and with the learner warm-started with the same knowledge (Section~\ref{sec:granularity}).
\end{enumerate}

\subsection{The cross-fitted headroom estimator}
\label{sec:estimator}
Let $Y_a(x)$ be the relative makespan of one run of action $a$ on instance $x$, with mean $\mu_a(x)$ and independent noise of variance $\sigma^2$; the headroom is $H(x)=\mu_{\mathrm{rand}}(x)-\min_a\mu_a(x)$. The \emph{one-seed oracle} $\hat H_1=\mu_{\mathrm{rand}}-\min_aY_a$ takes the best of one run per action; the \emph{cross-fitted oracle} chooses $\hat a=\arg\min_a\bar Y_a^{(1)}$ on two runs and evaluates it on two independent ones, $\hat H_{\mathrm{cf}}=\mu_{\mathrm{rand}}-\bar Y_{\hat a}^{(2)}$.
\begin{proposition}
\label{prop:bounds}
For every $x$, $\mathbb E[\hat H_{\mathrm{cf}}(x)]\le H(x)\le\mathbb E[\hat H_1(x)]$. If all $\mu_a(x)$ are equal, $\mathbb E[\hat H_{\mathrm{cf}}]=H=0$ and $\mathbb E[\hat H_1]\approx\sigma\sqrt{2\ln A}$ under normal noise.
\end{proposition}
The expectation of a minimum does not exceed the minimum of the expectations, so the one-seed value is an upper bound; conditional on $\hat a$, which depends only on the first batch, the second batch is unbiased for $\mu_{\hat a}\ge\min_a\mu_a$, so the cross-fitted value is a lower bound (proof and Monte Carlo check in the supplementary material). Every headroom statement below uses the cross-fitted value.

\subsection{Statistics, registration and reproducibility}
\label{sec:stats}
Per instance (the mean of the seeds), $d=100(A-B)/B$ is tested in pairs (Wilcoxon, two-sided), with Holm correction over each family, a bootstrap $90\%$ interval and an equivalence margin of $\pm0.05\%$~\cite{Lakens2017}: a Holm-significant difference whose interval lies inside the margin is \emph{parity}, and one whose interval crosses it a \emph{small win}; ``beyond the margin'' means confirmed and not equivalent to zero. The registered tests are pooled under a global Benjamini--Yekutieli control (Section~\ref{sec:discussion}). Hypotheses, decision rules and probabilistic predictions were written, with md5 hashes of the code and of the frozen inputs, in a dated, version-controlled research diary before each battery was run~\cite{NosekEtal2018}; each analysis script prints whether each prediction held. The code is compiled (numba) and covered by $1{,}577$ unit tests, and audit scripts recompute every number quoted in the text from the result files; all runs used one laptop (8 cores), where $K$ was measured.
